\documentclass[12pt, a4paper]{article}

% Paquetes requeridos para el circuito
\usepackage[utf8]{inputenc}
\usepackage[T1]{fontenc}
\usepackage{circuitikz} 

\begin{document}

\begin{center}
    \begin{circuitikz}[american, scale=1.0, transform shape]
        
        % 1. Raíl inferior común que representa el nodo 'b'
        \draw (0,0) -- (8,0) node[circ] {} node[right] {b};
        
        % 2. Rama Izquierda: Fuente independiente V1 entre 'b' y 'c'
        \draw (0,0) to[V, l=$V_1$] (0,4) node[circ] {} node[above] {c};
        
        % 3. Sección de Entrada: R1 horizontal, nodo 'd' y R2 vertical
        % Vx se mantiene en R2 como variable de control
        \draw (0,4) to[R, l=$R_1$] (3,4) node[circ] {} node[above] {d}
                    to[R, l=$R_2$, v=$V_x$] (3,0);
        
        % 4. Rama Central-Derecha: Fuente dependiente de corriente gm*Vx
        % Usamos 'cI' para el símbolo de fuente dependiente (rombo)
        \draw (5.5,4) to[cI, l=$g_m V_x$] (5.5,0);
        
        % 5. Rama Derecha: Resistencia R3 en paralelo
        \draw (8,4) to[R, l=$R_3$] (8,0);
        
        % Raíl superior derecho (formando el nodo 'a')
        \draw (5.5,4) -- (8,4) node[circ] {} node[above] {a};
        
    \end{circuitikz}
\end{center}

\end{document}